\chapter{I cardiofrequenzimetri e il controllo dell'allenamento}
\label{cap:cardiofrequenzimetri}

La rilevazione della \textbf{frequenza cardiaca} (FC) in telemetria è il modo più diffuso per
misurare l'\textbf{intensità} dello sforzo, e quindi il \textbf{carico interno}, direttamente sul
campo. Il capitolo parte dai requisiti dello strumento e dal comportamento della FC, passa alla
determinazione della $FC_{max}$ --- il riferimento su cui poggiano tutte le percentuali --- e
raccoglie poi i metodi con cui il carico interno viene quantificato: percezione dello sforzo,
metodo Edwards, TRIMP, metodo Lucia e prelievi di lattato.

% ---------------------------------------------------------------------------
\section{Innovazioni tecnologiche ed evoluzione}

\rubrica{Requisiti dei mezzi di valutazione dell'intensità}
\begin{itemize}
  \item \textbf{scientificamente attendibili};
  \item \textbf{non invasivi};
  \item di \textbf{facile somministrazione};
  \item \textbf{specifici}.
\end{itemize}

La rilevazione della FC rispetta tutti i parametri di \textbf{efficienza scientifica e
strumentale}, e rappresenta quindi un metodo utile per valutare l'intensità dello sforzo.

\rubrica{Che cosa consente la telemetria}
\begin{itemize}
  \item analizzare in \textbf{tempo reale} l'intensità di lavoro negli atleti;
  \item l'analisi diretta delle \textbf{risposte acute} permette di:
    \begin{itemize}
      \item \textbf{oggettivare} gli allenamenti;
      \item \textbf{personalizzarli};
      \item \textbf{modificarli}, anche in tempo reale;
      \item \textbf{ottimizzare} i tempi di recupero.
    \end{itemize}
\end{itemize}

% ---------------------------------------------------------------------------
\section{Il carico di allenamento e la classificazione dell'intensità}

\textbf{Carico di allenamento}: prodotto del \textbf{volume} per l'\textbf{intensità}.

\rubrica{Scala generale (qualitativa)}
È la classificazione dell'intensità più usata dai tecnici, e a volte anche dai preparatori:
\begin{enumerate}
  \item leggera;
  \item media;
  \item forte;
  \item intensa;
  \item massimale.
\end{enumerate}
Il limite evidente del metodo è la \textbf{soggettività}.

\rubrica{Classificazione in percentuale di un parametro conosciuto}
Metodo alternativo: esprimere l'intensità come percentuale di un parametro misurato. Nel caso
dell'allenamento della \textbf{resistenza} i parametri di riferimento sono:
\begin{itemize}
  \item la \textbf{frequenza cardiaca massima} ($FC_{max}$);
  \item la \textbf{velocità aerobica massimale} (VAM);
  \item la \textbf{velocità di soglia anaerobica} ($V_{obla}$);
  \item la \textbf{velocità record} sulla distanza proposta.
\end{itemize}

% ---------------------------------------------------------------------------
\section{Il comportamento della frequenza cardiaca}

\rubrica{FC media o distribuzione della FC?}
\begin{itemize}
  \item molte valutazioni si basano sulla \textbf{FC media di lavoro};
  \item la letteratura scientifica internazionale suggerisce invece di guardare alla
    \textbf{distribuzione} della FC nel corso della sessione allenante;
  \item i dati sono cioè riferiti alla \textbf{percentuale del tempo totale} dell'allenamento in cui
    l'atleta rimane nella zona (p.\,es.\ 85/90\%) della propria $FC_{max}$ --- che è anche la zona di
    riferimento dell'analisi delle partite.
\end{itemize}

\textbf{Perché la media non basta} (Sassi, 2005): in un'esercitazione ad alta intensità con la palla
quattro giocatori mostrano una \textbf{FC media praticamente uguale} (attorno all'80\% della
$FC_{max}$), ma una \textbf{distribuzione molto diversa} --- fra il giocatore ``c'' e il giocatore
``d'' lo scarto è del \textbf{16,8\%}. Considerando solo la FC media si sarebbero fatte riflessioni
non corrette al momento di verificare il lavoro effettivamente svolto e di programmare le sedute
successive.

\rubrica{La FC si adegua in ritardo al lavoro e alla pausa}
\begin{itemize}
  \item all'\textbf{inizio} del lavoro la FC diventa \textbf{stabile dopo 60$''$};
  \item alla \textbf{fine} del lavoro resta \textbf{stabile per 20$''$} prima di scendere;
  \item conseguenza: il valore letto sul cardiofrequenzimetro \textbf{non è istantaneo}, e in una
    esercitazione intermittente breve la FC non ha il tempo di rappresentare l'impegno reale.
\end{itemize}

\rubrica{Lavoro intermittente e lavoro continuo}
Confronto fra due esercitazioni condotte entrambe \textbf{alla VAM} e lette in \% della $FC_{max}$:
\begin{itemize}
  \item \textbf{intermittente 20$''$/20$''$ a navetta}: salita rapida fino a $\simeq 90\%$, poi
    andamento oscillante e \textbf{caduta anticipata} nella fase finale;
  \item \textbf{continuo in linea}: salita più graduale, \textbf{plateau} attorno al 90\% mantenuto
    più a lungo e discesa più tardiva.
\end{itemize}

% ---------------------------------------------------------------------------
\section{La determinazione della frequenza cardiaca massima}

Tutte le percentuali di lavoro poggiano sulla $FC_{max}$: il modo in cui la si determina condiziona
quindi l'intera prescrizione del carico.

\rubrica{Formule teoriche}
\begin{enumerate}
  \item \textbf{Cooper}: $FC_{max} = 220 - \textrm{et\`a}$;
    \begin{itemize}
      \item studi recenti hanno dimostrato una \textbf{variabilità interindividuale} della
        $FC_{max}$ fino al \textbf{10/15\%};
    \end{itemize}
  \item \textbf{Ball State University}:
    \begin{itemize}
      \item uomini: $FC_{max} = 214 - 0{,}8 \times \textrm{et\`a}$;
      \item donne: $FC_{max} = 209 - 0{,}7 \times \textrm{et\`a}$;
    \end{itemize}
  \item \textbf{Martin} (1989): $FC_{allenante} = FC_{max} - (0{,}45 \times FC_{riposo})$;
  \item \textbf{Cerretelli}: $FC_{max} = 216 - (1{,}1 \times \textrm{et\`a})$;
  \item \textbf{Tanaka}: $FC_{max} = 208 - 0{,}7 \times \textrm{et\`a}$.
\end{enumerate}

\rubrica{La frequenza cardiaca di riserva (Karvonen)}
\begin{itemize}
  \item \textbf{HRR} (\textit{heart rate reserve}): $HRR = FC_{max} - FC_{riposo}$; p.\,es.\
    $190 - 60 = 130$;
  \item \textbf{FC allenante}: $[220 - \textrm{et\`a} - HRR \times \% \textrm{ intensit\`a
    prescelta}] + HRR$;
  \item esempio per un atleta di 25 anni: $(195 - 130) \times 85\% + 130 = 185{,}25$.
\end{itemize}

\rubrica{Cooper e Tanaka a confronto (atleta di 25 anni)}
\begin{table}[htbp]
  \centering
  \small
  \renewcommand{\arraystretch}{1.3}
  \begin{tabularx}{\textwidth}{|X|c|c|}
    \hline
    \textbf{Carico interno FC} & \textbf{Cooper} ($220 - \textrm{et\`a}$) &
      \textbf{Tanaka} ($208 - 0{,}7 \times \textrm{et\`a}$) \\
    \hline
    80\% $FC_{max}$ & 156 & 152 \\
    \hline
    85\% $FC_{max}$ & 165,7 & 161,9 \\
    \hline
    90\% $FC_{max}$ & 175,5 & 171,45 \\
    \hline
  \end{tabularx}
\end{table}

Lo scarto fra le due formule è di \textbf{4--5 battiti} a ogni percentuale: applicato alla
prescrizione, sposta la zona di lavoro.

% ---------------------------------------------------------------------------
\section{Tre modi di rilevare la \texorpdfstring{$FC_{max}$}{FCmax}: lo studio di Spedicato e coll.}

\rubrica{Campione}
Quarantatré giovani calciatori: età $16{,}5 \pm 1{,}2$ anni; peso $70{,}0 \pm 7{,}5$\,kg; altezza
$177 \pm 6$\,cm; BMI $22{,}2 \pm 1{,}4$\,kg/m$^2$.

\rubrica{I tre metodi confrontati}
\begin{enumerate}
  \item \textbf{$FC_{max}$ in allenamento}: valore massimo di FC raggiunto in \textbf{cinque
    rilevamenti consecutivi}, effettuati in altrettante sedute durante la stagione regolare;
  \item \textbf{$FC_{max}$ nel test di Léger}: valore rilevato all'\textbf{ultimo step di corsa} che
    il giocatore è stato in grado di eseguire;
  \item \textbf{$FC_{max}$ da formula teorica}: \textbf{Cooper} ($220 - \textrm{et\`a}$) e
    \textbf{Tanaka} ($208 - 0{,}7 \times \textrm{et\`a}$).
\end{enumerate}

\rubrica{Risultati}
\begin{itemize}
  \item FC media rilevata: \textbf{allenamento} $195 \pm 13$\,bpm; \textbf{Léger} $193 \pm 10$\,bpm;
    \textbf{Cooper} $204 \pm 1$\,bpm; \textbf{Tanaka} $197 \pm 1$\,bpm;
  \item differenza \textbf{significativa} fra la $FC_{max}$ da formula di \textbf{Cooper} e le altre
    tre modalità ($p<0{,}001$);
  \item \textbf{nessuna} differenza significativa ($p>0{,}05$) fra allenamento, test di Léger e
    formula di Tanaka: le tre modalità risultano \textbf{statisticamente sovrapponibili};
  \item le due formule hanno una \textbf{dispersione minima} ($\pm 1$\,bpm) perché dipendono solo
    dall'età, mentre allenamento e Léger mostrano l'\textbf{ampia variabilità individuale} reale.
\end{itemize}

\rubrica{Conclusioni}
\begin{itemize}
  \item considerare la $FC_{max}$ per controllare l'allenamento è \textbf{necessario}, ma è
    altrettanto indispensabile determinarla con la \textbf{massima precisione} e con test che diano
    la misura reale di ciò che si vuole valutare;
  \item una scelta sbagliata del metodo porta a un'\textbf{errata considerazione dei carichi di
    lavoro} e a possibili insuccessi nella programmazione;
  \item la $FC_{max}$ rilevata \textbf{in allenamento} vale quanto quella da test di Léger e da
    formula di Tanaka, con il vantaggio di essere rilevata in condizioni \textbf{sport-specifiche};
  \item \textbf{ma} non è un test \textbf{riproducibile} con parametri standardizzati: è solo un
    mezzo di valutazione della $FC_{max}$ e \textbf{non dà possibilità comparative} in momenti
    diversi della stagione. Il confronto nel tempo richiede \textbf{test standardizzati}.
\end{itemize}

% ---------------------------------------------------------------------------
\section{I metodi di controllo dell'allenamento}

\begin{enumerate}
  \item \textbf{scala di Borg}:
    \begin{itemize}
      \item \textbf{RPE} (\textit{ratings perceived exertion}, valutazione dello sforzo percepito):
        \textbf{20 livelli};
      \item \textbf{CR10}: \textbf{10 livelli}, è quella più utilizzata;
    \end{itemize}
  \item metodo \textbf{Edwards Training Load} (5 zone di FC);
  \item metodo \textbf{Banister} (\textbf{TRIMP}, \textit{training impulse});
  \item metodo \textbf{Lucia};
  \item \textbf{prelievi di lattato ematico}.
\end{enumerate}

% ---------------------------------------------------------------------------
\section{La scala di Borg CR10 e il carico interno percepito}

\rubrica{I livelli della CR10}
Dal nullo allo sforzo massimo: 0 nullo; 0,5 estremamente debole; 1 molto debole; 2 debole;
3 moderato; 5 forte; 7 molto forte; 10 estremamente forte; ``?'' massimo assoluto (Borg, 1981, 1982,
1998).

\rubrica{Corrispondenza con la \texorpdfstring{$FC_{max}$}{FCmax}}
\begin{itemize}
  \item livelli \textbf{5 e 6} $\rightarrow$ \textbf{85\%} $FC_{max}$;
  \item livelli \textbf{7 e 8} $\rightarrow$ \textbf{90/95\%} $FC_{max}$.
\end{itemize}

\rubrica{Gli indici ricavabili dalla CR10}
\begin{enumerate}
  \item \textbf{carico di allenamento} (\textit{training load}):
    \[ TL = \textrm{durata UA (min)} \times RPE \]
    espresso in \textbf{UA} (unità arbitrarie); p.\,es.\ $110 \times 4 = 440$\,UA;
  \item \textbf{indice di monotonia}: descrive la situazione in cui il TL di diverse unità di
    allenamento è simile;
    \[ IM = \frac{\textrm{media carico settimanale}}{\textrm{deviazione standard settimanale}} \]
  \item \textbf{fatica acuta} o \textit{strain}:
    \[ FA = TL \times IM \]
    tanto più il carico di allenamento è stato \textbf{elevato e monotono}, tanto più la fatica acuta
    sarà elevata.
\end{enumerate}

\rubrica{Valori di riferimento nel calcio}
\begin{itemize}
  \item \textbf{singola seduta}: l'effetto allenante si ottiene per valori di \textit{session} RPE
    fra \textbf{300 e 600\,UA};
  \item \textbf{somma settimanale}: fra \textbf{2000 e 3000\,UA};
  \item soglia di lettura: $>2600$\,UA carico \textbf{alto}, $<2600$\,UA carico \textbf{basso}.
\end{itemize}

\rubrica{Training load di una partita}
\begin{itemize}
  \item \textbf{alto}: $>700$\,UA;
  \item \textbf{medio-alto}: 700--500\,UA;
  \item \textbf{medio}: 500--400\,UA;
  \item \textbf{medio-basso}: 400--300\,UA;
  \item \textbf{basso}: $<300$\,UA.
\end{itemize}

\rubrica{Percezione degli allenatori e percezione degli atleti}
Confronto del \textit{training load} attribuito alle stesse esercitazioni da allenatori e atleti
(Foster e coll., 2001):
\begin{itemize}
  \item esercitazioni \textbf{\textit{easy}}: gli atleti percepiscono un carico \textbf{più alto} di
    quello attribuito dagli allenatori (differenza significativa);
  \item esercitazioni \textbf{\textit{moderate}}: le due percezioni \textbf{coincidono}
    sostanzialmente;
  \item esercitazioni \textbf{\textit{hard}}: gli allenatori percepiscono un carico \textbf{più alto}
    di quello riferito dagli atleti (differenza significativa);
  \item conseguenza pratica: il carico \textbf{va chiesto all'atleta}, non stimato dallo staff.
\end{itemize}

% ---------------------------------------------------------------------------
\section{Il metodo Edwards (1993)}

\textbf{Principio}: si misura la \textbf{durata} dell'allenamento (in minuti) trascorsa in
\textbf{5 diverse zone di FC}, ciascuna moltiplicata per il proprio \textbf{coefficiente}, e si
\textbf{sommano} i risultati.

\rubrica{Le cinque zone}
\begin{itemize}
  \item zona 1: 50--60\% $FC_{max}$ $\rightarrow$ coefficiente 1 ($K_1$);
  \item zona 2: 60--70\% $FC_{max}$ $\rightarrow$ 2 ($K_2$);
  \item zona 3: 70--80\% $FC_{max}$ $\rightarrow$ 3 ($K_3$);
  \item zona 4: 80--90\% $FC_{max}$ $\rightarrow$ 4 ($K_4$);
  \item zona 5: 90--100\% $FC_{max}$ $\rightarrow$ 5 ($K_5$).
\end{itemize}

\rubrica{Edwards modificato: la ``zona partita''}
Nel monitoraggio della FC in situazioni di gioco e in esercitazioni tecnico-tattiche condizionali le
bande sono espresse in \% della $HR_{peak}$:
\begin{itemize}
  \item $HR_{peak}$ 90--100\%, 85--90\% e 80--95\% costituiscono la \textbf{zona partita};
  \item $HR_{peak}$ 75--80\%: banda intermedia;
  \item $HR_{peak}$ 70--75\%: \textbf{recupero};
  \item $HR_{peak}$ $<70\%$: \textbf{riposo}.
\end{itemize}

\textbf{Percentuale di lavoro}: quantità di tempo, in percentuale, in cui l'atleta è in
\textbf{zona partita}. Le esercitazioni sono \textbf{intense} quando la \% di lavoro è
$>70\%$, cioè quando per almeno il 70\% del tempo di esercitazione il giocatore è stato in zona
partita.

% ---------------------------------------------------------------------------
\section{Il metodo TRIMP di Banister (1986)}

\textbf{TRIMP} (\textit{training impulse}): indice di carico interno calcolato come
\[ TRIMP = TD \cdot \Delta FC \cdot y \]
dove:
\begin{itemize}
  \item \textbf{TD}: durata totale dell'allenamento;
  \item \textbf{$\Delta FC$}: frequenza cardiaca di riserva,
    \[ \Delta FC = \frac{FC_{media\ esercizio} - FC_{riposo}}{FC_{max} - FC_{riposo}} \]
  \item \textbf{y}: coefficiente \textbf{non lineare}, $y = b\,e^{cx}$ ovvero
    $y = 0{,}64\,e^{1{,}92x}$, con $e$ base del logaritmo neperiano e $x = \Delta FC$;
    $b = 0{,}64$ e $c = 1{,}92$ \textbf{per i maschi}.
\end{itemize}

\rubrica{Il limite del metodo}
Banister utilizza la \textbf{FC media} dell'esercizio di tutta la sessione e un fattore $y$
calcolato con due costanti ($b$, $c$) che sono \textbf{uguali per tutti i soggetti}.

% ---------------------------------------------------------------------------
\section{Il TRIMP individuale di Manzi}

\rubrica{Il problema}
La FC media dell'esercizio e uno \textbf{stesso} fattore $y$ di moltiplicazione non riescono a
esprimere le \textbf{richieste fisiologiche individuali} di ogni sessione di allenamento.

\rubrica{La soluzione}
\begin{itemize}
  \item si introduce per \textbf{ogni soggetto} un \textbf{fattore individuale di ponderazione}
    ($y_i$), che riflette il profilo della sua tipica \textbf{curva di risposta del lattato ematico}
    all'aumentare dell'intensità;
  \item il fattore si ricava dalla \textbf{linea esponenziale} che interpola i punti
    lattato/elevazione frazionale della FC del singolo atleta --- p.\,es.\
    $y_i = 0{,}2445\,e^{3{,}411x}$ con $R^2 = 0{,}97$;
  \item all'aumentare dell'intensità (letta dalla risposta della FC registrata) il fattore $y_i$
    aumenta in modo \textbf{esponenziale};
  \item il \textbf{$TRIMP_i$} può quindi essere calcolato in \textbf{qualsiasi momento} della
    sessione allenante e quantificato come lo \textbf{pseudo-integrale} di tutti i punti della curva
    $\Delta FC$.
\end{itemize}

% ---------------------------------------------------------------------------
\section{La distribuzione dell'intensità e la fitness aerobica (Manzi)}

\rubrica{Impianto dello studio}
\begin{itemize}
  \item l'utilità di quantificare il TL con i metodi della FC si valuta mettendo in relazione:
    \begin{enumerate}
      \item l'\textbf{intensità di allenamento} (tempo passato in una determinata zona di intensità);
      \item i \textbf{risultati} dell'allenamento (variabili della \textbf{fitness aerobica
        submassimale}, come le soglie del lattato);
    \end{enumerate}
  \item \textbf{scopo}: esaminare la distribuzione dell'intensità, quantificata con la FC, durante la
    \textbf{prima fase di preparazione pre-campionato} (6 settimane), periodo in cui le squadre di
    Serie A svolgono un rilevante programma di condizionamento fisico;
  \item l'intensità è stata quantificata con \textbf{3 zone di FC} riferite ai risultati del
    \textbf{test submassimale al treadmill}:
    \begin{itemize}
      \item \textbf{bassa} intensità: FC $<2$\,mmol/l;
      \item \textbf{moderata} intensità: FC fra 2 e 4\,mmol/l;
      \item \textbf{alta} intensità: FC $>4$\,mmol/l.
    \end{itemize}
\end{itemize}

\rubrica{Risultati (504 allenamenti individuali)}
\begin{itemize}
  \item distribuzione del tempo di allenamento:
    \begin{itemize}
      \item bassa intensità: $73 \pm 2{,}5\%$;
      \item media intensità: $19 \pm 2{,}8\%$;
      \item alta intensità: $8 \pm 1{,}4\%$;
    \end{itemize}
  \item FC alle soglie: a 2\,mmol/l $81{,}5 \pm 2{,}6\%$ della $FC_{max}$; a 4\,mmol/l
    $90{,}2 \pm 2{,}1\%$;
  \item velocità alla soglia di 2\,mmol/l: $10{,}9 \pm 0{,}6$\,km/h pre e $11{,}4 \pm 0{,}5$\,km/h
    post allenamento;
  \item velocità alla soglia di 4\,mmol/l: $13{,}0 \pm 0{,}8$\,km/h pre e $13{,}9 \pm 0{,}4$\,km/h
    post allenamento.
\end{itemize}

\rubrica{Effetti pratici nel calcio}
\begin{itemize}
  \item la fitness aerobica submassimale è influenzata \textbf{positivamente} dalla quantità di
    esercizio ad \textbf{alta intensità} sostenuta in allenamento;
  \item nel calciatore d'élite va allenata con esercizi che prevedano fasi ad alta intensità per
    \textbf{almeno l'8\%} del programma totale di allenamento settimanale;
  \item quindi: per migliorare la fitness aerobica dei calciatori servono \textbf{esercitazioni ad
    alta intensità}.
\end{itemize}

% ---------------------------------------------------------------------------
\section{Il metodo Lucia}

\begin{itemize}
  \item il \textbf{Lucia's TRIMP} è di \textbf{difficile applicazione}: richiede \textbf{test di
    laboratorio} per determinare i parametri dell'allenamento;
  \item \textbf{calcolo}: si moltiplica il tempo (in minuti) speso in \textbf{tre} zone per il
    coefficiente $K$ relativo a ogni zona e si sommano i risultati:
    \begin{itemize}
      \item zona 1: \textbf{sotto la soglia ventilatoria} $\times K_1$;
      \item zona 2: fra la \textbf{soglia ventilatoria} e il \textbf{punto di recupero
        respiratorio} $\times K_2$;
      \item zona 3: \textbf{al di sopra} del punto di recupero respiratorio $\times K_3$;
    \end{itemize}
  \item è un metodo \textbf{simile a quello di Edwards}, ma basato su \textbf{parametri di
    laboratorio} anziché su percentuali della $FC_{max}$.
\end{itemize}

% ---------------------------------------------------------------------------
\section{I prelievi di lattato ematico}

\rubrica{Riferimenti}
\begin{itemize}
  \item \textbf{2\,mmol/l}: \textbf{soglia aerobica};
  \item \textbf{4\,mmol/l}: \textbf{soglia anaerobica};
  \item p.\,es.\ per lavorare \textbf{sopra soglia} in un lavoro di \textbf{potenza aerobica}
    l'atleta deve trovarsi al di sopra delle 4\,mmol/l.
\end{itemize}

\rubrica{Lattato, FC e sforzo percepito}
\begin{itemize}
  \item la \textbf{combinazione} di FC e concentrazione di lattato ematico predice con
    \textbf{discreta precisione} la RPE;
  \item all'aumentare di FC e lattato aumenta anche lo \textbf{sforzo percepito}.
\end{itemize}

\rubrica{Perché si misura il lattato e non l'acido lattico}
\begin{itemize}
  \item il \textbf{lattato} è un \textbf{sale} o \textbf{estere} dell'acido lattico;
  \item l'\textbf{acido lattico} vero e proprio è difficile da misurare: si trova in ambiente
    acquoso all'interno della fibra muscolare ed è \textbf{dissociato in due ioni} --- lo
    \textbf{ione lattato} con carica negativa e lo \textbf{ione idrogeno} con carica positiva.
\end{itemize}

% ---------------------------------------------------------------------------
\section{Esempi dal campo}

I tracciati di FC di atleti che svolgono la \textbf{stessa} seduta --- esercitazioni intermittenti
sui 50--100\,m, $5 \times 50$\,m in 7$''$ per 3 serie, circuiti fisico-coordinativi (allenamenti
della Dinamo Bucarest, stagione 2011/12) --- mostrano \textbf{risposte individuali molto diverse}:

\begin{itemize}
  \item alcuni atleti restano in zona di FC \textbf{molto elevata} per quasi tutta la durata della
    sessione;
  \item altri si mantengono su valori \textbf{più bassi}, con innalzamenti sopra i 170\,bpm solo
    \textbf{verso la fine} del lavoro o solo \textbf{nell'ultima serie} delle tre proposte;
  \item nello stesso circuito, un atleta mostra un andamento \textbf{soddisfacente} con soli piccoli
    picchi oltre l'intensità massima, un altro si attesta \textbf{quasi sempre} sull'intensità
    massima sopportata.
\end{itemize}

\rubrica{Come si interpretano}
\begin{itemize}
  \item carico \textbf{sovrastimato} per l'atleta con FC costantemente elevata;
  \item carico \textbf{sottostimato} per l'atleta con FC bassa;
  \item oppure, semplicemente, \textbf{stato di allenamento diverso} fra i due;
  \item in ogni caso la stessa proposta non produce lo stesso carico interno: il
    \textbf{cardiofrequenzimetro} è un mezzo valido per determinare il carico di allenamento
    \textbf{individuale} del singolo calciatore.
\end{itemize}
